The architecture of prominent LLMs, which is primarily based on the Transformer model \cite{vaswani2017attention}, has enabled a significant increase in scale and performance of language models. Further, the development from foundational models, which were primarily trained on next-token prediction, such as GPT 3, to instruction-tuned and chat-optimized models, like GPT 3.5 and GPT 4, has also marked a significant shift. Techniques like Reinforcement Learning from Human Feedback (RLHF) have been used to align these models more closely with user intent, making them more adept at following instructions and engaging in coherent, multi-turn conversations \cite{ouyang2022training}. This project specifically evaluates these newer, chat-optimized models, as their interactive nature makes them more fitting candidates for playing dialogue games.

The \inlinecode{clembench} framework \cite{chalamalasetti-etal-2023-clembench,beyer2024clembench} serves as the foundation for this project. Its purpose is to enable the systematic and reproducible evaluation of LLMs as conversational agents through game play. The framework is structured such that a programmatic game master controls the game flow, presents instructions and game states to LLM players, and parses their responses. This allows for automated scoring and isolates the evaluation to the abilities of the LLMs to understand and act within the constraints of the game, rather than their ability to generate unconstrained text.

Existing \inlinecode{clemgames} in the benchmark evaluate a wide range of different abilities. For instance, the taboo game tests concept description and semantic understanding under constraints, while the drawing and reference games test a form of multimodal grounding through character-based grids. The DoND game presented in this project complements these by introducing a direct test of negotiation, compromise, and strategic reasoning based on conflicting objectives and incomplete information. These capabilities are central to many real-world interactions but were not the primary focus of prior games included in the benchmark. Other related work in evaluating LLMs in interactive environments also includes frameworks that use LLMs as referees and focus on more open-ended social interactions \cite{li2023beyond,zhou2024sotopia}. In contrast, \inlinecode{clembench} focuses on structured, automatically scorable games, providing a possibly more easily understandable and objective lens.

The problem of dividing resources between agents with different preferences is a classic topic in game theory, and has been formalized by John Nash in 1950 \cite{nash1950equilibrium,nash1950bargaining}. The DoND game is such a multi-issue bargaining task, where agents must find a point in the solution space that is mutually agreeable. Note however that unlike in the standard semi-competitive setting, in the cooperative setting also evaluated in this project, the task is purely cooperative, and so the agreement should be trivial for rational agents if given full information.

Research on AI agents in negotiation also has a long history, though much of it has focused on agents that communicate using symbolic actions or pre-defined logical forms rather than unrestricted natural language. However, a number of key papers have used the DoND setting to train AI systems. This project is substantially influenced by those works. \cite{lewis2017deal} and \cite{he2018decoupling} used this and similar negotiation tasks to train agents using a combination of supervised and reinforcement learning. More recently, \cite{liao2024efficacy} used the same DoND game to investigate the efficacy of fine-tuning through self-play. They found that iterative training via filtered behaviour cloning could dramatically improve model performance. However, they concluded that the increased performance was not due to improved high-level strategy, but due to improvements in rule-following and the ability to reach agreements.

This project fundamentally differs from these works in its objective. The goal here is not to train or fine-tune models. Instead, the DoND game is used as an instrument for zero-shot evaluation of pre-trained, general-purpose LLMs, assessing their inherent out-of-the-box negotiation capabilities. This way the project benchmarks the baseline performance of these models and identifies their intrinsic strengths and weaknesses in a negotiation context.
